% =====================================================================
\part{A apresentação}
% =====================================================================

\chapter{Roteiro de vinte minutos}

O princípio que organiza este roteiro: \textbf{a apresentação não é um resumo do
documento}. Ela é o argumento, e o argumento tem uma forma --- quatro perguntas
encadeadas, em que cada resposta obrigou a seguinte. Uma apresentação que
percorre capítulos na ordem do sumário gasta metade do tempo em metodologia e
chega aos resultados com a banca já cansada.

Duas regras práticas. \textbf{Um slide, uma afirmação.} E \textbf{toda
afirmação-manchete vem com o seu limite no mesmo slide}, porque limite dito
por você é rigor, e limite descoberto pela banca é falha.

\vspace{10pt}
\begin{longtable}{@{}p{0.9cm}p{4.5cm}p{9.2cm}@{}}
\toprule
\textbf{Min.} & \textbf{Slide} & \textbf{O que dizer} \\
\midrule
\endfirsthead
\toprule
\textbf{Min.} & \textbf{Slide} & \textbf{O que dizer} \\
\midrule
\endhead
0--1 & Título e a frase & A tese em uma frase. Nada mais neste slide. \\
\addlinespace
1--3 & O problema & Goiás já entra ocupado em 1985. O que as quatro décadas
registram é reorganização, não abertura de fronteira nova. Três famílias de
explicação concorrem, e elas \emph{recomendam políticas opostas}. \\
\addlinespace
3--4 & Objetivos & Os quatro objetivos específicos, que são as quatro frentes.
Anuncie que a apresentação segue essa ordem. \\
\addlinespace
4--7 & Dados e unidades & MapBiomas 10.1 + IBGE + Ipeadata + SICOR. Os 246
municípios que viram 166 AMCs, e \emph{por quê} (as duas fontes se comportam de
maneiras opostas diante de uma emancipação). Primeira diferença em tudo, e a
pergunta que ela reformula. \\
\addlinespace
7--9 & Os três atos & Figura do painel de cobertura em cinco cortes. As duas
fronteiras e a âncora externa: a soja do IBGE quebra em 2020 sozinha. Diga que
o Ato~I é o mais destrutivo. \\
\addlinespace
9--11 & \textbf{Frente 1} --- o padrão & Figura dos centros de massa. Os quatro
números com IC. \textbf{Diga ``ancorada''} para a vegetação natural. Cite as
duas verificações mais fortes: pixel a pixel e os controles negativos (urbana
para o sul; leite menos da metade). \\
\addlinespace
11--13 & \textbf{Frente 2} --- o mecanismo & Figura da idade da pastagem. Duas
lógicas simultâneas; a assimetria de $-49\%$ contra $-13\%$. O contraste 45\%
contra 1\% de lavoura anterior. E o achado negativo: a região explica 0,5\%. \\
\addlinespace
13--16 & \textbf{Frente 3} --- o empurrão & Figura das doze especificações. As
duas assinaturas procuradas; as doze negativas; o poder do teste temporal. O
que ocupa o lugar: substituição local. Depois, o motor comum e o veredito
\emph{corroborante, não estabelecido}. \\
\addlinespace
16--18 & \textbf{Frente 4} --- o teto & Figura da decomposição. A demanda no
pico; o Sul acelerando 51\% em pasto e caindo pela metade em Cerrado; a taxa
que cai com o esgotamento. Os 6,56 Mha, ditos como \emph{exposição}. \\
\addlinespace
18--19 & O que fica e o que não fica & O quadro dos graus de estabelecimento,
inteiro, na tela. Este slide é o seu escudo: ele mostra que você já graduou a
própria evidência. \\
\addlinespace
19--20 & Cronograma & Revisão de literatura como frente mais extensa; as três
arestas com desenho definido; depósito em dezembro, defesa em janeiro. \\
\bottomrule
\end{longtable}

\section{Os cinco momentos em que você ganha a banca}

\begin{enumerate}
\item \textbf{Ao dizer ``ancorada''} em vez de ``andou 7,6 km''. Explique em uma
frase por quê. É a primeira demonstração de que você conhece a diferença entre
sinal e ruído.
\item \textbf{Ao mostrar que o erro de medida aponta contra a tese.} Os 46 km ao
norte do centro da massa reetiquetada. É contraintuitivo e memorável.
\item \textbf{Ao apresentar o nulo do empurrão com o poder declarado} antes de
alguém perguntar.
\item \textbf{Ao rebaixar você mesmo a Frente 3b} de significante para
corroborante, dizendo que o $p$ agrupado era otimista.
\item \textbf{Ao dizer que a primeira regressão do teto de oferta não é
resultado nenhum}, e sim uma identidade contábil se reapresentando.
\end{enumerate}

\section{Os quatro erros que custam caro}

\begin{enumerate}
\item Gastar mais de sete minutos antes da primeira frente. Metodologia é
importante e está inteira no documento; a apresentação precisa chegar ao
argumento.
\item Dizer ``refutei'' em qualquer contexto.
\item Apresentar a união de Agricultura com Mosaico como ``a medida correta''.
\item Deixar as autocorreções para responder na arguição. Elas são o seu maior
ativo e devem aparecer \emph{na sua fala}, ainda que em um slide só.
\end{enumerate}

\begin{nabanca}[Se a apresentação tiver de encolher para 15 minutos]
Corte os três atos para um slide só, funda dados e unidades num único slide, e
preserve intactas as quatro frentes e o quadro dos graus. O que nunca se corta é
o slide dos graus de estabelecimento.
\end{nabanca}

\chapter{Cartão de bolso}

% Os blocos abaixo são pseudotítulos, e não seções — não devem entrar no
% sumário. Sem o \Needspace, um deles cai no pé da página com o conteúdo na
% seguinte, que num cartão de consulta é o defeito que mais custa: procura-se
% pelo rótulo.
\newcommand{\bloco}[1]{%
  \par\Needspace*{5\baselineskip}\vspace{10pt}%
  {\color{cortinta}\large\bfseries #1}\par\vspace{2pt}}

Os números que precisam sair da boca sem consulta.

\bloco{O estado e o dado}

\vspace{2pt}
246 municípios $\to$ \textbf{166 AMCs} (53 grupos + 113 idênticas) \quad$|$\quad
Painel $246 \times 40 = 9.840$ linhas \quad$|$\quad \textbf{38} diferenças
anuais \quad$|$\quad MapBiomas Coleção 10.1, 30 m, 1985--2024 \quad$|$\quad
\textbf{58} rotinas, \textbf{31} decisões, \textbf{38} referências

\bloco{O balanço de 40 anos}
Vegetação natural \textbf{17,65 $\to$ 11,88 Mha} ($51{,}9\% \to 34{,}9\%$)
\quad$|$\quad Agricultura \textbf{1,17 $\to$ 5,58 Mha} ($\times 4{,}8$)
\quad$|$\quad Soja $\times 12$ em área, $\times 13$ em produção
\quad$|$\quad Pastagem \textbf{10,98 $\to$ 11,99 Mha}, com pico $\approx 14{,}8$
em \textbf{2003} \quad$|$\quad Rebanho $+46\%$ (15,9 $\to$ 23,2 mi de cabeças)
com área de pasto $+9\%$

\bloco{Os três fluxos}
Veg. natural $\to$ pastagem \textbf{4,10 Mha} \quad$|$\quad Pastagem $\to$
agricultura \textbf{2,72 Mha} \quad$|$\quad Veg. natural $\to$ agricultura
\textbf{1,29 Mha} \quad$|$\quad A pastagem é o elo de \textbf{58\%} da entrada
em agricultura

\bloco{A periodização}
\textbf{2001} ($F = 62{,}2$) e \textbf{2020} ($F = 21{,}5$) \quad$|$\quad Atos:
\textbf{1985--2000}, \textbf{2001--2019}, \textbf{2020--2024} \quad$|$\quad
Âncora externa: a soja do IBGE quebra em 2020 sozinha

\bloco{Frente 1 --- o padrão}
Pastagem \textbf{$+77{,}6$ km} $[54{,}7; 98{,}2]$ \quad$|$\quad Rebanho
\textbf{$+66{,}9$} $[47{,}2; 84{,}5]$ \quad$|$\quad Agricultura
\textbf{$+65{,}2$} $[43{,}5; 94{,}6]$ \quad$|$\quad Vegetação natural
$+7{,}6$ $[-0{,}5; 15{,}6]$ = \textbf{ancorada} \\
Vão lavoura--pasto: \textbf{123 a 135 km} em todos os 40 anos \quad$|$\quad
Pixel a pixel confirma dentro de 2 km \quad$|$\quad Urbana anda \textbf{8 km ao
sul}

\bloco{Frente 2 --- o mecanismo}
Veg.$\to$pasto cai \textbf{49\% no Sul} e \textbf{13\% no Norte} \quad$|$\quad
\textbf{44,6 mi} de conversões pasto$\to$lavoura (3,8 Mha), \textbf{16,0 mi} com
idade conhecida \\
Duas populações: $\mu \approx 4$ anos (33\%) e $\mu \approx 16$ anos (67\%)
\quad$|$\quad Origem lavoura: mediana \textbf{5 anos}; origem vegetação:
\textbf{13 anos} \\
Até 3 anos: \textbf{45\%} vieram de lavoura; 30+ anos: \textbf{1\%}
\quad$|$\quad A mesorregião explica \textbf{0,5\%}

\bloco{Frente 3 --- empurrão e motor comum}
\textbf{12 especificações espaciais, todas negativas} (a hipótese exigia
positivas) \quad$|$\quad \textbf{24 combinações temporais}, $p$ mínimo
\textbf{0,078} \\
Poder: \textbf{93\%} contra empurrão forte, \textbf{48\%} contra moderado
\quad$|$\quad Substituição local: $\mathbf{-0{,}52}$ a $\mathbf{-1{,}14}$ \\
\emph{Shift-share}: $\beta = \mathbf{-0{,}033}$, $p$ de permutação
\textbf{0,07--0,13} \quad$|$\quad Piso da régua: $1/38 = 0{,}026$ \\
Corrida: a aptidão perde \textbf{62\%} da magnitude com a latitude no modelo
\quad$|$\quad $r(\text{aptidão}, \text{latitude}) = -0{,}44$

\bloco{Frente 4 --- o teto de oferta}
Fluxo bruto estadual \textbf{0,071 $\to$ 0,072 Mha/ano} (o endereço mudou, não a
velocidade) \quad$|$\quad Sul: $-37\%$ \\
No Sul do Ato III: pasto$\to$lavoura \textbf{$+51\%$}, expansão da soja
\textbf{$+244\%$}, Cerrado novo cai \textbf{pela metade} \\
Demanda no pico: câmbio $134{,}5 \to 169{,}0$; soja $104{,}4 \to 186{,}4$
\quad$|$\quad Decomposição: \textbf{17\%} estoque, \textbf{83\%} residual \\
Taxa cai \textbf{0,5 a 0,8 p.p.} por décimo de esgotamento --- negativo em
\textbf{16 de 16} tratamentos, cruza 5\% em \textbf{11} \\
\textbf{6,56 Mha} ainda \emph{expostos} ($\approx 60\%$ de 1985), Norte concentra
\textbf{44\%} \quad$|$\quad \textbf{6,35} deles fora de Proteção Integral \\
Remanescente florestal do Sul: \textbf{52,2\% $\to$ 59,9\%} (três quintos)

\bloco{Carbono}
Estoque removido \textbf{973 Mt} CO\textsubscript{2}e \quad$|$\quad Emissão
líquida \textbf{833 Mt} (desconto de 14,4\%) \\
Savânica \textbf{573 Mt} contra florestal \textbf{340 Mt} --- é resultado sobre
\emph{área}, não sobre densidade \\
1 ha de floresta = \textbf{1,6} ha de savânica = \textbf{2,6} ha de campestre
\quad$|$\quad Três quartos (\textbf{722 Mt}) saíram no \textbf{Ato I}

\bloco{As quatro palavras que não podem escapar}
\textbf{Ancorada} (não ``andou'') \quad$\cdot$\quad \textbf{Exposta} (não
``disponível'') \quad$\cdot$\quad \textbf{Corroborante} (não ``significativo'')
\quad$\cdot$\quad \textbf{Removido} (não ``emitido'', nas contas por ato)

\vspace{14pt}
{\color{corregra}\rule{\textwidth}{0.8pt}}

\vspace{6pt}
{\color{cortinta}\bfseries E a frase para fechar a defesa:}

\vspace{2pt}
\emph{``Goiás não abriu uma fronteira nova entre 1985 e 2024: reorganizou a que
já tinha. E a diferença entre essas duas leituras é a diferença entre uma
política que realoca atividade e uma que age sobre o retorno de converter ---
porque a primeira redistribui a conversão, e não a reduz.''}
